\documentclass[11pt,a4paper]{article}
\usepackage[margin=23mm]{geometry}
\usepackage{amsmath,amssymb}
\usepackage{xcolor}
\usepackage{hyperref}
\definecolor{accent}{HTML}{175C78}
\hypersetup{colorlinks=true,urlcolor=accent,linkcolor=accent,pdftitle={Bellman equations: a practical guide for Wall-E RL},pdfauthor={Wall-E RL project}}
\setlength{\parindent}{0pt}
\setlength{\parskip}{7pt}
\newcommand{\block}[1]{\vspace{8pt}{\large\bfseries\color{accent}#1}\par}
\newcommand{\note}[1]{\begin{quote}\color{accent}\textbf{#1}\end{quote}}
\begin{document}
{\small\color{accent} WALL-E RL\quad /\quad LEARNING REFERENCE\hfill 27 September 2026}\par
\vspace{12pt}

{\Huge\bfseries Bellman equations}\par
{\Large Reward now, value later}

\note{The value of a situation equals the expected reward from the next action, plus the discounted value of the situation that follows.}

This guide connects Bellman equations to our CartPole experiment and PPO's critic. Work through the numbers first; the notation describes the same idea. All numerical examples below are illustrative, not outputs from our saved model.

\block{1. Reward, return, and value}
CartPole gives $+1$ for each step, including the failure step in our default setup. But a nearly upright pole and a rapidly falling pole can have the same immediate reward and very different futures.

\begin{tabular}{p{25mm}p{122mm}}
\textbf{Reward} & The number received from one transition.\\[5pt]
\textbf{Return} & The accumulated rewards along an attempt, possibly discounted.\\[5pt]
\textbf{Value} & The expected return from a starting situation under a specified policy.
\end{tabular}

The critic predicts value. It does not choose the action; that is the actor's job.

\block{2. A three-step example}
Suppose exactly three rewards remain before a true terminal state:
\[
s_0\ \xrightarrow{\ r=1\ }\ s_1\ \xrightarrow{\ r=1\ }\ s_2\ \xrightarrow{\ r=1\ }\ \text{terminal}.
\]
With no discounting, the return is $1+1+1=3$. With discount $\gamma=0.9$:
\[
G_0=1+0.9(1)+0.9^2(1)=2.71.
\]
From the next state, two rewards remain, so $G_1=1+0.9(1)=1.9$. Therefore:
\[
\boxed{G_0=1+0.9G_1=1+0.9(1.9)=2.71.}
\]
The remaining future has been summarized by one number. This decomposition is the starting point for Bellman's relationship.

\block{3. Why it works}
Let $r_t$ be the reward received after the action at step $t$. Then
\begin{align*}
G_t &= r_t+\gamma r_{t+1}+\gamma^2r_{t+2}+\cdots\\
&=r_t+\gamma\bigl(r_{t+1}+\gamma r_{t+2}+\cdots\bigr)\\
&=\boxed{r_t+\gamma G_{t+1}}.
\end{align*}
Bellman equations apply this decomposition to \emph{expected} returns.

\newpage
{\small\color{accent} BELLMAN EQUATIONS\hfill POLICY EVALUATION}\par
\block{4. The value belongs to a policy}
A policy $\pi$ is a rule for choosing actions. Starting from the same CartPole state, an always-right controller and a trained controller can have different expected returns.
\[
V^\pi(s)=\text{expected discounted return from state $s$, following $\pi$}.
\]
The superscript $\pi$ identifies the policy; it is not an exponent. Its Bellman equation is
\[
\boxed{V^\pi(s)=\mathbb{E}\!\left[r_t+\gamma V^\pi(s_{t+1})\mid s_t=s\right].}
\]
\begin{tabular}{p{25mm}p{122mm}}
$s$ & Current state.\\[4pt]
$s_{t+1}$ & Next state after the action and environment update.\\[4pt]
$r_t$ & Reward received from that transition.\\[4pt]
$\gamma$ & Discount factor; later rewards get smaller weights.\\[4pt]
$\mathbb{E}$ & Average over policy action choices and possible environment outcomes.
\end{tabular}

For this equation, the state must contain enough information to predict the next-state distribution from the current state and action. A partial sensor observation may need history or an estimator to provide that information. Do not assume Wall-E's pitch alone is a complete state.

\block{5. Connect it to a stochastic policy}
In a toy state, suppose the policy chooses left with probability $0.3$ and right with probability $0.7$. Each action leads to one definite next state in this example:

\begin{center}
\begin{tabular}{lrrr}
\textbf{Action} & \textbf{Probability} & \textbf{Reward} & \textbf{Next value}\\\hline
Left & 0.3 & 1 & 10\\
Right & 0.7 & 1 & 20
\end{tabular}
\end{center}
With $\gamma=0.9$, the reward-plus-next-value is $1+0.9(10)=10$ for left and $1+0.9(20)=19$ for right. Thus
\[
\boxed{V^\pi(s)=0.3(10)+0.7(19)=16.3.}
\]
The average follows the policy's probabilities. One sampled action does not itself represent this entire expectation.

\block{6. State value versus action value}
$V^\pi(s)$ lets the policy choose the first action. $Q^\pi(s,a)$ fixes that first action to $a$, then follows $\pi$ afterward:
\[
Q^\pi(s,a)=\mathbb{E}\!\left[r_t+\gamma V^\pi(s_{t+1})\mid s_t=s,a_t=a\right].
\]
For the toy state, $Q^\pi(s,\text{left})=10$ and $Q^\pi(s,\text{right})=19$. Their policy-weighted average gives $V^\pi(s)=16.3$:
\[
V^\pi(s)=\sum_a\pi(a\mid s)Q^\pi(s,a).
\]

\newpage
{\small\color{accent} BELLMAN EQUATIONS\hfill OPTIMALITY AND LEARNING}\par
\block{7. Evaluate a policy or find the best one?}
The expectation equation asks how well a specified policy performs. The optimal value $V^*(s)$ asks for the best expected return achievable from that state:
\[
\boxed{V^*(s)=\max_a\mathbb{E}\!\left[r_t+\gamma V^*(s_{t+1})\mid s_t=s,a_t=a\right].}
\]
\textbf{Policy evaluation:} average action outcomes using the policy's probabilities.\par
\textbf{Optimality:} select the action with the greatest expected outcome, assuming optimal behavior afterward too.

An optimal action-value function obeys
\[
Q^*(s,a)=\mathbb{E}\!\left[r_t+\gamma\max_{a'}Q^*(s_{t+1},a')\mid s_t=s,a_t=a\right].
\]
The maximum is over actions, not over lucky environment outcomes. The latter still require an expectation. PPO's critic approximates the current policy's value; it is not directly told the optimal value.

\block{8. Learn from a one-step prediction error}
Let $\hat V$ denote a critic's approximate value. Suppose a transition has
\[
\hat V(s_t)=15,\qquad r_t=1,\qquad \hat V(s_{t+1})=20,\qquad\gamma=0.9.
\]
A one-step target is
\[
y_t=r_t+\gamma\hat V(s_{t+1})=1+0.9(20)=19.
\]
The \textbf{temporal-difference error} compares that target with the old prediction:
\[
\boxed{\delta_t=r_t+\gamma\hat V(s_{t+1})-\hat V(s_t)=4.}
\]
A table-based update with learning rate $\alpha=0.1$ would give
\[
V_{\mathrm{new}}(s_t)=15+0.1(4)=15.4.
\]
A neural-network critic instead changes its weights to reduce prediction errors over many samples. Using an estimated next value in the target is \textbf{bootstrapping}. It avoids waiting for the whole future, but an inaccurate next prediction can make the target inaccurate too.

\block{9. The end of an episode}
After a true terminal transition, no future reward remains. Use $y_t=r_t$. The default CartPole failure step earns $1$, so its one-step target is $1$, even though the attempt failed.

Do not use the reset state's value as the continuation of the failed attempt. A time-limit truncation can be different: if it only interrupts an otherwise continuing task, the next value can still matter. If the deadline is part of the task itself, value depends on remaining time. This is why \texttt{terminated} and \texttt{truncated} need different handling.

\newpage
{\small\color{accent} BELLMAN EQUATIONS\hfill PPO AND QUICK REFERENCE}\par
\small
\block{10. Where this appears in our PPO experiment}
PPO's Generalized Advantage Estimation (GAE) combines temporal-difference errors across steps. Away from episode boundaries, its pattern is
\[
\hat A_t\approx\delta_t+(\gamma\lambda)\delta_{t+1}
+(\gamma\lambda)^2\delta_{t+2}+\cdots.
\]
The rollout length and episode boundaries determine where this sum ends. Our configuration uses $\gamma=0.99$ and $\lambda=0.95$.

\begin{center}
\fbox{\begin{minipage}{0.88\linewidth}
Collect transitions and critic predictions\par
$\downarrow$\par
Use rewards and next values to estimate advantages and return targets\par
$\downarrow$\par
Train the critic toward the targets; train the actor using advantages\par
$\downarrow$\par
Collect fresh experience with the updated policy
\end{minipage}}
\end{center}

Positive estimated advantage favors increasing an action's probability; negative advantage favors decreasing it. PPO's clipping modifies the actor objective to reduce incentives for excessive policy changes. The critic uses rollout-based targets, not only the simple one-step target illustrated on page 3.

\block{11. Why a score of 500 is not a value of 500}
Our reported episode score sums undiscounted rewards. If exactly 500 rewards of $+1$ remain, the discounted sum at $\gamma=0.99$ is
\[
\sum_{k=0}^{499}0.99^k=\frac{1-0.99^{500}}{1-0.99}\approx99.34.
\]
The episode score is still 500. The critic's output is an estimate, so it need not equal this exact hypothetical sum. Neither number establishes physical robot performance.

\block{12. Keep these distinctions clear}
\begin{itemize}
\setlength{\itemsep}{2pt}
\item Reward is one transition; value summarizes the expected future.
\item Bellman is a consistency relationship, not a neural-network architecture.
\item A stochastic policy requires averaging across its possible choices.
\item A positive TD error is a prediction surprise, not proof of causation.
\item Bellman relationships do not require knowing the physics equations: samples can provide reward and next-state information.
\end{itemize}

\block{References and project files}
Primary reference: OpenAI Spinning Up, \href{https://spinningup.openai.com/en/latest/spinningup/rl_intro.html}{Part 1: Key Concepts in RL}, sections on value functions, Bellman equations, and advantages. The \href{https://github.com/openai/spinningup/blob/master/docs/spinningup/rl_intro.rst}{official source text} was checked for this guide.

Project connections:\\
\texttt{configs/ppo.yaml} --- experiment settings.\\
\texttt{training/train.py} --- learning and checkpoints.\\
\texttt{training/common.py} --- environment and evaluation.\\
\texttt{docs/CARTPOLE\_GYMNASIUM.md} --- runnable walkthrough.
\par The actor runs at inference; the critic and optimizer support laptop training.
\end{document}
